\chapter{The DSL's static guarantees, in detail}
\label{ch:dsl-guarantees}

\begin{aside}
Chapter~\ref{ch:dsl} introduced the DSL's type-state machine.
This chapter zooms in on the specific guarantees: what the DSL
prevents at compile time, what it allows at runtime, and how
each guarantee maps to a constraint or concept. Treat it as a
reference for ``what can I assume about my code if it
compiled?''.
\end{aside}

\section{Guarantee 1: borderless boxes cannot have border colour}

\begin{lstlisting}
auto a = box(text(\"x\"));                       // no border
auto b = a | bcol<60, 65, 80>;                  // compile error
\end{lstlisting}

Mechanism: the \code{operator|(BoxNode<Cfg>, bcol\_t)} has a
\code{requires (Cfg.has\_border)} clause. Without a border, the
clause fails.

Why it matters: a border colour without a border produces no
visible effect at runtime --- a silent bug. The compile error
makes the intent explicit.

\section{Guarantee 2: padding is non-negative}

\begin{lstlisting}
auto a = text(\"x\") | pad<-1>;                  // compile error
\end{lstlisting}

Mechanism: the \code{pad} template uses an NTTP constrained to
non-negative integers via a concept.

Why it matters: negative padding has no defined meaning; a
runtime crash or visual glitch would be the alternative.

\section{Guarantee 3: each pipe attribute applies at most once}

\begin{lstlisting}
auto a = text(\"x\") | width<10> | width<20>;    // compile error
\end{lstlisting}

Mechanism: the second \code{| width<20>} has a \code{requires
(!Cfg.has\_width)} clause; the first pipe set the flag, so the
second fails.

Why it matters: ``which width wins?'' would be ambiguous;
forbidding the double-pipe makes the intent explicit. If you
really want to override, build a fresh node.

\section{Guarantee 4: children must be elements}

\begin{lstlisting}
auto a = h(text(\"x\"), 42);                     // compile error
\end{lstlisting}

Mechanism: \code{h(\ldots)} requires all arguments to satisfy
\code{ElementLike}. An \code{int} does not.

Why it matters: silent conversion would produce a runtime crash
or empty element.

\section{Guarantee 5: NTTPs are constant expressions}

\begin{lstlisting}
std::string s = compute();
auto a = t<s>;                                 // compile error
\end{lstlisting}

Mechanism: \code{t<S>} takes a \code{Str} NTTP; only literals
and constexpr values qualify.

Why it matters: the compile-time path is exclusive to compile-
time strings. For runtime strings, use \code{text(s)} explicitly.

\section{Guarantee 6: style composition is associative}

\begin{lstlisting}
auto a = Bold | Italic | Underline;
auto b = (Bold | Italic) | Underline;
auto c = Bold | (Italic | Underline);
// a, b, c are the same type.
\end{lstlisting}

Mechanism: \code{StyTag<A> | StyTag<B>} returns \code{StyTag<merge(A,
B)>}, and \code{merge} is associative on the underlying flag
struct.

Why it matters: refactoring style expressions is safe; the
order of pipes does not affect the result.

\section{Guarantee 7: TextNode preserves literal storage}

\begin{lstlisting}
auto a = t<\"hello\">;                           // TextNode<Str{\"hello\"}>
// At runtime, the bytes \"hello\" are in the static string table.
// The runtime element holds a string_view into them.
\end{lstlisting}

Mechanism: the NTTP carries the literal; the runtime element
holds a \code{string\_view} that points into the static segment.

Why it matters: no allocation for compile-time text; no
copy-on-construction.

\section{Guarantee 8: Style merge precedence is deterministic}

\begin{lstlisting}
auto a = text(\"x\") | Bold;                     // Bold
auto b = box(a) | Italic;                       // Bold + Italic
// The order of Bold and Italic across the parent/child is unambiguous.
\end{lstlisting}

Mechanism: the merge function is pure and deterministic; same
inputs, same output.

Why it matters: behavioural reasoning is local; reading a pipe
chain tells you the resulting style without running the program.

\section{Guarantee 9: width-clipped text never exceeds width}

\begin{lstlisting}
auto a = text(very_long) | width<10>;
// At runtime, the rendered text is exactly 10 cells wide.
\end{lstlisting}

Mechanism: the layout phase respects the width as a hard cap.
The text renderer truncates (with optional ellipsis) when
content exceeds.

Why it matters: layout stability; widgets that allocate 10
cells get exactly 10.

\section{Guarantee 10: flex direction is one-way}

\begin{lstlisting}
auto a = h(text(\"x\")) | direction<Column>;     // ?? not allowed
\end{lstlisting}

Mechanism: \code{h(\ldots)} creates a row; the direction is
baked in. Overriding via pipe is not supported; if you want a
column, write \code{v(\ldots)}.

Why it matters: ambiguity in direction would force a runtime
choice; making it explicit at the DSL level avoids the
confusion.

\section{Guarantees the DSL does not give}

For completeness, the things the DSL does \emph{not} guarantee:

\begin{itemize}[leftmargin=*]
  \item Runtime content fits the layout. A long runtime string
    might overflow its width container.
  \item Terminal capabilities. The DSL emits truecolour even
    if the terminal does not support it; quantisation happens
    later.
  \item Sub-tree count. A vector of children can be of any
    length; the DSL does not constrain.
  \item Memory budget. A deeply nested tree can blow the
    stack; the DSL does not bound recursion.
\end{itemize}

These are runtime concerns that the DSL leaves to the
runtime. The discipline is: use the DSL for shape, use the
runtime for content.

\section{Cumulative effect}

The ten guarantees, plus a dozen more that fall out of concept
constraints, mean a \maya{} UI that compiles is structurally
sound. You cannot ship a panel with a border colour and no
border; you cannot pass an \code{int} where an element is
expected; you cannot accidentally set width twice.

This is what the type-state machine buys you: the compiler
caught the bugs before you ran the program.

\begin{idea}
A compiler that catches structural bugs is the best code review
you will ever get. The DSL turns the C++ compiler into a UI
linter, automatically and continuously. Every successful build is
a small certificate of correctness.
\end{idea}

\section{Mapping guarantees to error messages}

When a guarantee fires, the compiler emits an error. The shape
is consistent:

\begin{lstlisting}
error: no match for 'operator|' (operand types are X and Y)
note: candidate: 'requires Z ...'
note:   constraint not satisfied: 'Z evaluated to false'
\end{lstlisting}

Read top-down: ``no match'' is the call site; ``candidate'' is
the overload that almost matched; ``constraint not satisfied''
is the specific rule. Map the rule to one of the guarantees in
this chapter; the fix follows.

\section{Guarantees in widget code}

Widget builders can declare their own constraints. If a widget
requires certain configuration, encode it as a concept:

\begin{lstlisting}
template <class T>
concept HasLabel = requires (T t) { { t.label } -> std::convertible_to<std::string_view>; };

template <HasLabel T>
Element labelled(const T& t, std::string_view extra);
\end{lstlisting}

Now callers without a \code{label} field get a clear error,
not an SFINAE wall.

\section*{Workshop}

\begin{enumerate}[leftmargin=*]
  \item Pick three of the ten guarantees. For each, write code
    that violates it, read the resulting error, identify which
    \code{requires} clause caught it.
  \item Write a small widget with a concept-constrained
    parameter. Verify the error message is clear when the
    concept is violated.
  \item Read the source of a DSL pipe overload. Confirm the
    \code{requires} clause is human-readable.
\end{enumerate}

\begin{recap}
The DSL gives ten structural guarantees, all enforced at compile
time. They eliminate whole classes of bugs (silent decorations,
ambiguous widths, wrong child types). Errors map back to a
specific \code{requires} clause; the fix is mechanical. The
trade-off is some DSL complexity; the payoff is a UI that
compiles is a UI that is structurally sound.
\end{recap}
